% Part: many-valued-logic
% Chapter: infinite-valued-logics
% Section: goedel

\documentclass[../../../include/open-logic-section]{subfiles}

\begin{document}

\olfileid{mvl}{inf}{god}

\olsection{منطقات غودل}

\begin{editorial}
  ما يأتي نواة موجزة لقسم يُبسط فيه القول في منطق غودل اللانهائي القيم.
\end{editorial}

\begin{defn}\ollabel{def:goedel} يعرّف منطق غودل اللانهائي القيم
$\LogGod[\infty]$ بالمصفوفة الآتية:
\begin{enumerate}
  \item لغة القضايا القياسية~$\Lang L_0$ بروابطها
  $\lfalse$ و$\lnot$ و$\land$ و$\lor$ و$\lif$.
  \item مجموعة قيم الصدق~$V_\infty$.
  \item $1$ هي القيمة المميزة الوحيدة؛ أي~$V^+ = \{1\}$.
  \item تعطي الدوال الآتية دوال الصدق:
  \begin{align*}
    \tf{\lfalse} & = 0\\
    \tf{\lnot}[\LogGod](x) & = \begin{cases}
      1 & \text{إذا كان } x =0\\
      0 & \text{في غير ذلك}
    \end{cases}\\
    \tf{\land}[\LogGod](x,y) & = \min(x,y)\\
    \tf{\lor}[\LogGod](x,y) & = \max(x,y)\\
    \tf{\lif}[\LogGod](x,y) & = \begin{cases}
      1 & \text{إذا كان } x \le y\\
      y & \text{في غير ذلك.}
    \end{cases}
    \end{align*}
\end{enumerate}
ويعرّف منطق غودل ذو~$m$ من القيم بالطريقة نفسها، عدا أن~$V = V_m$.
\end{defn}

\begin{prop}
  المنطق~$\LogGod[3]$ المعرّف في \olref[thr][god]{defn:goedel} هو
  نفسه~$\LogGod[3]$ المعرّف في \olref{def:goedel}.
\end{prop}

\begin{proof}
  يتضح ذلك بمقارنة جداول صدق الروابط المعطاة في
  \olref[thr][god]{defn:goedel} بجداول الصدق التي تحددها المعادلات في
  \olref{def:goedel}:
  \begin{center}
    \begin{tabular}{c|c}
      $\tf{\lnot}[\LogGod[3]]$ & \\
      \hline
      $1$ & $0$ \\
      $1/2$ & $0$ \\
      $0$ & $1$
    \end{tabular}
    \quad
    \begin{tabular}{c|ccc}
      $\tf{\land}[\LogGod]$ & $1$ & $1/2$ & $0$ \\
      \hline
      $1$ & $1$ & $1/2$ & $0$ \\
      $1/2$ & $1/2$ & $1/2$ & $0$\\
      $0$ & $0$ & $0$ & $0$
    \end{tabular}
    \\[2ex]
    \begin{tabular}{c|ccc}
      $\tf{\lor}[\LogGod]$ & $1$ & $1/2$ & $0$ \\
      \hline
      $1$ & $1$ & $1$ & $1$ \\
      $1/2$ & $1$ & $1/2$ & $1/2$ \\
      $0$ & $1$ & $1/2$ & $0$
    \end{tabular}
    \quad
    \begin{tabular}{c|ccc}
      $\tf{\lif}[\LogGod]$ & $1$ & $1/2$ & $0$ \\
      \hline
      $1$ & $1$ & $1/2$ & $0$ \\
      $1/2$ & $1$ & $1$ & $0$  \\
      $0$ & $1$ & $1$ & $1$
    \end{tabular}
  \end{center}
\end{proof}

\begin{prop}\ollabel{prop:god-infty-m}
  إذا كان~$\Gamma \Entails[\LogGod[\infty]] !B$، فإن
  $\Gamma \Entails[\LogGod[m]] !B$ لكل~$m \ge 2$.
\end{prop}

\begin{proof}
  تمرين.
\end{proof}

\begin{prob}
  أثبت \olref[mvl][inf][god]{prop:god-infty-m}.
\end{prob}

% تصحيح مصدر (0401-I1): لا يصح العكس المعلن هنا إلا مع مقدمات متناهية.
وإذا كانت مجموعة المقدمات متناهية، فليس اللزوم مقصورًا على هذا الاتجاه؛
بل يصح عكسه أيضًا.

وكما كان الأمر في~$\LogGod[3]$، فإن~$\LogGod[\infty]$ يجعل كل~!!{formula}s الصحيحة
حدسيًا تحصيلات منطقية. وأما الأمثلة التي سبق بيان أنها ليست تحصيلات،
فتبقى كذلك في~$\LogGod[\infty]$:
\begin{align*}
  & p \lor \lnot p && (p \lif q) \lif (\lnot p \lor q) \\
  & \lnot\lnot p \lif p && \lnot(\lnot p \land \lnot q) \lif (p \lor q) \\
  & ((p \lif q) \lif p) \lif p && \lnot(p \lif q) \lif (p \land \lnot q)
\end{align*}
وأما مثال~!!{formula} التي لا تصح حدسيًا مع كونها تحصيلًا منطقيًا
في~$\LogGod[3]$، أعني~$(p \lif q) \lor (q \lif p)$، فهو تحصيل منطقي
في~$\LogGod[\infty]$ أيضًا. ولنا في~$\LogGod[\infty]$ وصف آخر:
نأخذ المنطق الحدسي ونضيف إليه المخطط~$(!A \lif !B) \lor (!B \lif !A)$.
وهذا الوصف أثبته مايكل دَمِت؛ ومن هنا شاع أن يسمى~$\LogGod[\infty]$
منطق غودل--دَمِت، ورمزه~$\Log{LC}$.

\begin{prob}
  بين أن~$(p \lif q) \lor (q \lif p)$ تحصيل منطقي في
  $\LogGod[\infty]$.
\end{prob}

\begin{prob}
  بين أن~$(p \lif q) \lor (q \lif r) \lor (r \lif s)$، وهو تحصيل
  منطقي في~$\LogGod[3]$، ليس تحصيلًا منطقيًا في~$\LogGod[\infty]$.
\end{prob}

\end{document}
